\chapter{Seguridad y modelo de amenazas}
\label{cap:seguridad-amenazas}

\section{Activos, adversarios y fronteras}
El modelo de amenazas del proyecto aplica STRIDE a las interfaces, servicios,
almacenamientos y canales de correo. Los activos principales son la identidad,
la participación, las credenciales del votante, la urna, la cadena de eventos,
los resultados y los datos de cada organización. Se consideran adversarios
externos, votantes maliciosos, gestores, usuarios de otros tenants, servicios
comprometidos, operadores y observadores de red. El detalle y la matriz de
riesgos se encuentran en \texttt{docs/architecture/threat-model.md}.

\section{Controles implementados}
Los controles documentados incluyen tokens almacenados como hashes, límites
de intentos para códigos, tokens de un solo uso, aislamiento por
\texttt{organization\_id}, validación del cliente de servicio esperado en
rutas internas, segregación de las bases de urna y censo, registros
administrativos de solo adición y verificación de la cadena antes del
escrutinio. La seguridad externa presupone TLS y que solo el gateway esté
expuesto.

% [Bloque F] Delegación de credenciales frente a las prácticas (RGPD).
Este planteamiento difiere del seguido en las Prácticas Externas. En
TaskFlow, la propia API almacenaba las contraseñas de los usuarios protegidas
con bcrypt, de modo que la aplicación asumía su custodia. En VeriVoto, los
servicios propios no almacenan contraseñas: la autenticación de los usuarios
de gestión se delega en Keycloak, y el votante accede mediante un enlace
personal y un código de un solo uso, que se guardan como hashes, o mediante su
cuenta institucional. Se reduce así la cantidad de datos sensibles que debe
proteger la lógica de negocio, en línea con el principio de minimización del
RGPD. Del mismo modo que las claves VAPID en las prácticas, los secretos de
configuración se proporcionan mediante variables de entorno y no forman parte
del repositorio.

\section{Riesgos residuales y límites}
El correo es el canal de enlace y código para procesos configurados con
\texttt{EMAIL\_CODE}; quien controle ese buzón puede suplantar al votante. La
cuenta institucional reduce ese riesgo bajo los controles de D-48, pero no
elimina el requisito de proteger la identidad federada. Un operador con acceso
a las dos bases podría intentar correlaciones temporales o alterar ambas
fuentes. También permanecen fuera de alcance la resistencia a coacción, el
malware en el dispositivo, la disponibilidad ante ataques a gran escala y la
verificabilidad extremo a extremo. Por ello, no se afirma anonimato absoluto
frente al operador de la plataforma.
